\documentclass[11pt,letterpaper]{article}
\usepackage[T1]{fontenc}
\usepackage{lmodern,microtype,amsmath,amssymb,booktabs,array,enumitem}
\usepackage[letterpaper,left=1.02in,right=1.02in,top=.90in,bottom=.92in]{geometry}
\usepackage{xcolor,fancyhdr}
\usepackage[pdfencoding=auto,psdextra]{hyperref}
\definecolor{ink}{HTML}{173B49}
\definecolor{link}{HTML}{1B5968}
\hypersetup{colorlinks=true,linkcolor=link,citecolor=link,urlcolor=link,pdftitle={The 432 imaginary quadratic fields of class-group exponent six}}
\urlstyle{same}
\pagestyle{fancy}\fancyhf{}
\fancyhead[L]{\footnotesize\color{ink} Transfers from Dirichlet GRH}
\fancyhead[R]{\footnotesize\color{ink} Research note}
\fancyfoot[C]{\small\thepage}
\setlength{\headheight}{14.2pt}
\renewcommand{\headrulewidth}{.25pt}
\setlength{\emergencystretch}{3em}
\setlength{\parindent}{0pt}\setlength{\parskip}{.44em plus .08em minus .06em}
\setlist{topsep=.2em,itemsep=.15em,parsep=0pt}
\setcounter{tocdepth}{1}
\raggedbottom
\newcommand{\repo}[2]{\href{https://github.com/lsernau/grh-lean/blob/72fead102f663d5c7a407964a13dc0cf9d977207/grh/#1}{#2}}
\begin{document}
\thispagestyle{plain}
{\small\sffamily\bfseries\color{link} MATHEMATICAL RESEARCH NOTE\hfill 25 NOVEMBER 2026}\par
\vspace{.30in}
{\fontsize{27}{32}\selectfont\bfseries\color{ink} The 432 imaginary quadratic fields\\ of class-group exponent six\par}
\vspace{.11in}
{\large A consequence of Dirichlet GRH, with applications to prime-gap graphs and quintic reflection fields}\par
\vspace{.21in}

\textbf{Main theorem.} Exactly 432 imaginary quadratic fields have ideal class group of exponent six. Their largest absolute fundamental discriminant is $5\,761\,140$. Elsenhans, Kl\"uners and Nicolae established the list through $3.1\cdot10^{20}$. Their GRH bound left a possible interval extending to $2.5\cdot10^{25}$. We exclude the entire interval. It follows that the corresponding table of Aiken, Ross, Shvydkoy and Xue for quadratic orders is complete: it contains 1,236 discriminants.

\textbf{Why the bound closes.} A signed test in the explicit formula forces a hypothetical omitted field to have at least 28 split primes below 5500. Their norm equations admit at most 18 positive values for $b$ in $a^2+|D|b^2=4p^6$. A shared coefficient eliminates $D$ between two equations. Every one of the remaining $9\,393\,302$ candidate norm values has a small prime witness that contradicts the exponent-six assumption. The finite checks use integer and rational arithmetic.

\textbf{Other transfers.} The repository also discharges the RH hypothesis in a published proof of Toroczkai's two prime-gap graph conjectures. We give a connected growth history with permanent neighborhoods. For quintic fields we answer the incoming report's local discriminant question and exhibit a single prime-splitting test equivalent to RH for six nonabelian zeta functions. At exponent seven we obtain a smaller upper bound; completeness remains open.
\smallskip
\tableofcontents

\clearpage
\section{What the Lean development supplies}
The analysis uses commit \texttt{72fead102f663d5c7a407964a13dc0cf9d977207} of Luke Sernau's repository, the shorter development based on Levent Alp\"oge's work with Claude. The two editions pin Lean 4.35.0 and Mathlib \texttt{56a005855ea17e0b0958f73ec801ae140b0a1679}. We compared the public statement definitions, checked the five exports in both source trees and inspected the proof's last reduction. The included source-audit program checks these claims against a local checkout. Rubinstein's independent compact-edition verification records successful compilation, an axiom audit and kernel replay of 794 of 794 modules. We rely on the published replay for kernel verification. The downstream proofs here are human-readable, with separate finite checks.
\repo{src/Statements.lean}{Public statements};
\repo{src/Main.lean}{exported theorems};
\href{https://mrubinst.github.io/grh-lean-compact-check/}{independent compact-edition verification}.

\subsection*{The public propositions}
\texttt{riemannHypothesis} has Mathlib's type \texttt{RiemannHypothesis}. The second export spells out its nontrivial zeros.
\texttt{LFunction\_ne\_zero\_of\_half\_lt\_re} gives, for all $N\ge1$ and all complex Dirichlet characters $\chi$ modulo $N$,
\[
  \Re s>1/2,\quad (\chi\ne1\text{ or }s\ne1)\quad\Longrightarrow\quad L(s,\chi)\ne0.
\]
The other two exports, \texttt{LFunction\_eq\_zero\_re} and \texttt{LFunction\_eq\_zero\_re\_of\_isPrimitive}, give the zero location in the open critical strip. The reported axiom list consists of \texttt{propext}, \texttt{Classical.choice}, and \texttt{Quot.sound}. For an abelian number field over $\mathbb Q$, factorization of the Dedekind zeta function into primitive Dirichlet factors supplies the relevant field RH. In particular the quadratic zeta hypothesis in the published discriminant cutoff is available for every candidate field.

\paragraph{The part of the proof that transfers.}
The arithmetic construction places a weighted sum of prime powers in a comparison of covolumes. Short integral forms and a spanning estimate bound a covolume from above. Determinants of chains of theta lifts bound it from below. Their comparison yields the source inequality $m\,R\le f a(\log b)^5$ for the constructed remainder and parameters. After passage to the analytic remainder this produces a one-sided bound of order $x^{-1/2+\epsilon}$. The prime weights use a congruence condition and an auxiliary quadratic twist. A Mellin transform then detects zeros through poles whose real residues have the needed sign. Landau's theorem rules out those poles using the single side of the remainder bound. Section~\ref{sec:joint-test} transports this last step to a family whose splitting behavior can be written down explicitly. The geometric estimate in the source concerns its own congruence and twist weights.
\repo{GUIDE.md}{Reading guide};
\repo{src/New/Reduction/MainInequality.lean}{covolume comparison};
\repo{src/New/Reduction/OneSidedBound.lean}{analytic bound};
\repo{src/New/Reduction/Zeros/Criterion.lean}{zero criterion}.

\paragraph{Earlier work used here.}
At ProveIt commit \texttt{1e1728ad221d69e48fea09bf754b73422af44943}, \texttt{docs/incoming} contained three research archives. Their roles in this note can be read in the following order.
\begin{itemize}
\item \texttt{grh-lean-transfers.zip} lists the 26 quintic CM coefficient fields and asks for the exact local behavior at five. Its cubic example illustrates how the sign of a prime weight detects zeros.
\item \texttt{grh-transfer.zip} treats Zhao's spectral condition for 229 abelian fields and gives a completely split prime criterion for general normal extensions. Its criterion motivates comparing the splitting test for one field with a simultaneous one for six.
\item \texttt{grh\_transfers.zip} supplies the exponent-five base list and earlier restrictions on exponent-six discriminants. It proves completeness when at most six primes ramify or when the cofactor after removing the largest prime is at most $10^8$.
\end{itemize}
The computations below generate the radical polynomials from norm equations and recompute their maximal-order discriminants. The only imported finite field list for this part is the table of 27 quadratic inputs. We do not claim publication priority for the deductions or catalogs.
\href{https://github.com/VladimirReshetnikov/ProveIt/tree/1e1728ad221d69e48fea09bf754b73422af44943/docs/incoming}{Pinned directory of prior reports}.

\par\medskip\section{Completing the exponent-six census}
\label{sec:zero-test}
\textbf{Theorem.} The 432 fields in the exponent-six row of Elsenhans, Kl\"uners and Nicolae's Table~1 are all the imaginary quadratic number fields of exact class-group exponent six. Among them the largest absolute discriminant is $5\,761\,140$. A quadratic class group has exponent six when its sixth-power map is identically one and six is the least positive exponent with this property.

The finite enumeration of Elsenhans, Kl\"uners and Nicolae (EKN) is complete for $|D|\le X_0=3.1\cdot10^{20}$. Their Table~2, assuming the quadratic case of Dirichlet GRH, bounds the discriminant of every exponent-six field by $X_1=2.5\cdot10^{25}$. Indeed the analytic input is a bound for a rational prime splitting in $\mathbb Q(\sqrt D)$, whose zeta function is $\zeta(s)L(s,\chi_D)$. These are exactly the two zero statements available from the repository. We show that an exponent-six field with $X_0<|D|\le X_1$ cannot exist.
\href{https://arxiv.org/pdf/1803.02056}{Elsenhans, Kl\"uners and Nicolae (EKN), Theorems 1--3 and Tables 1--2};
\href{https://doi.org/10.1090/S0025-5718-96-00763-6}{Bach and Sorenson, Table~3}.

Write $x=|D|$, $e$ for the class-group exponent and
$\chi(n)=(D/n)$ for the primitive odd quadratic character of conductor $x$.
If $p$ splits and $P$ is either prime above $p$, then $P^e$ is principal.
Its generator $(a+b\sqrt D)/2$ has $b\ne0$ because its valuations at $P$
and $\bar P$ differ. Norms give $x\le 4p^e$. Consequently, if
$x>4Y^e$, every prime $p\le Y$ is inert or ramified. The explicit
formula can distinguish this pattern from the one required by the zeros
of $L(s,\chi)$.

\medskip\noindent\textbf{An explicit-formula inequality.}
Let $g$ be a compactly supported, real step function with nonzero norm.
Set
\[
 F(u)=\frac{\int_{\mathbb R}g(v)g(v+u)\,dv}{\int_{\mathbb R}g(v)^2\,dv},
 \qquad h(t)=\int_{\mathbb R} F(u)e^{itu}\,du.
\]
Then $F$ is even, continuous, piecewise linear, with $F(0)=1$ and
$|F(u)|\le1$. Its Fourier transform $h$ is nonnegative on the real line,
being a squared modulus there. It extends to an entire function with
$O((1+|t|)^{-2})$ decay on every fixed horizontal strip, by the explicit
Fourier integral of a step function and its reflected copy. Write $C$
for the largest absolute slope of $F$. For $\psi=\Gamma'/\Gamma$ the
Guinand--Weil formula is
\[
 \sum_{\rho}h\big((\rho-1/2)/i\big)=
 \log(x/\pi)+\frac{1}{2\pi}\int_{\mathbb R}h(t)
        \Re\psi(3/4+it/2)\,dt
 -2\sum_{n\ge2}\frac{\chi(n)\Lambda(n)}{\sqrt n}F(\log n).
\]
The left side is nonnegative by the repository's Dirichlet GRH. This is
where the zeros enter the bound. The odd parity determines the argument
$3/4$ of the digamma function.
\href{https://arxiv.org/pdf/2404.08380}{Carneiro, Milinovich,
Quesada-Herrera and Ramos, Lemma~9, with Fourier conventions converted}.

By the integral formula for the digamma function and Fourier inversion,
the integral including its factor $1/(2\pi)$ equals
\[
 \psi(3/4)+
 2\int_0^\infty \frac{e^{-3u/2}}{1-e^{-2u}}(1-F(u))\,du
 \ \le\ \psi(3/4)+\frac{14}9 C .
\]
Indeed the difference in the digamma integrand involves $1-\cos(tu)$;
its nonnegativity against $h(t)$ permits interchange of integrals.
Use $1-F(u)\le C u$, expand the denominator geometrically, and bound
$\sum_{j\ge0}(2j+3/2)^{-2}\le4/9+1/3$ by the integral test.
Also $\log\pi-\psi(3/4)>2$. To see this, $\log\pi>1.1$ and
$\gamma>H_6-\log 7>1/2$, while the first two terms of the digamma
difference give $\psi(1)-\psi(3/4)>1/3+1/14$ and $\psi(1)=-\gamma$.
It follows that
\[
 \log x> 2\sum_{n\ge2}\frac{\chi(n)\Lambda(n)}{\sqrt n}
                    F(\log n)+2-\frac{14}9 C.
\]

\medskip\noindent\textbf{A finite split-prime test.}
For each support bound $Y$ take $g$ constant on twelve successive intervals of width
$\delta=\log(Y)/12$, beginning at zero, with respective values
\[
          21,\ 14,\ 10,\ 6,\ 4,\ 1,\ -1,\ -4,\ -6,\ -10,\ -14,\ -21
\]
and zero elsewhere. Thus $F$ vanishes for $|u|\ge\log Y$. The opposite
signs at the ends allow inert primes to give a positive term in the
explicit formula even though $\chi(p)=-1$. Define
\[
 R_p=2\sum_{p^k\le Y}\frac{(-1)^k}{\sqrt{p^k}}F(k\log p),\qquad
 S=\sum_{p\le Y}R_p\log p,\qquad B=\max(0,\max_{p\le Y}R_p).
\]
If no $p\le Y$ splits, the prime sum including its factor two is
$S-\sum_{p\mid D,\,p\le Y}R_p\log p\ge S-B\log x$.
Hence that splitting pattern forces
\[
                 \log x>\frac{S+2-14C/9}{1+B}.
\]

\subsection{Forcing two equal coefficients}
Assume $e=6$ and $X_0<x\le X_1$. Because $4\cdot2000^6<X_0$, all split primes exceed 2000. Take $Y=5500$. Let $N_Y$ be the number of split primes up to $Y$. For these primes only $k=1$ occurs in the explicit formula. Changing a prime from inert to split subtracts the quantity
\[
   L_p=-4(\log p) F(\log p)/\sqrt p
\]
from the inert-model sum. The contribution lost at all ramified primes is at most $B\log x$, since $\sum_{p\mid D}\log p\le\log x$. Thus, for any common upper bound $L>0$ for the split-prime losses,
\[
 N_Y L > S+2-14C/9-(1+B)\log X_1.
\]
The exact rational enclosures in \texttt{zero\_test.py} give the two strict inequalities
\[
 \max_{2000<p\le5500} L_p<0.216,\qquad
 S+2-14C/9-(1+B)\log X_1>5.84.
\]
For orientation $S$ is about $65.682$, $B$ about $0.04476$, and $C$ about $0.47355$. Since $27\cdot0.216=5.832<5.84$, there are at least 28 split primes. Every such prime gives a norm equation
\[
       a_p^2+x b_p^2=4p^6,\qquad a_p\ge0,\quad b_p\ge1 .
\]
Here one may take absolute values of the two coefficients. The exact inequality $4\cdot5500^6<19^2 X_0$ gives $b_p\le18$. By the pigeonhole principle two distinct primes $p<q$ share the same $b$.

\subsection{Eliminating the discriminant}
Write $A=a_p$, $B'=a_q$. These nonnegative integers have the same parity, since their squares differ by a multiple of four. Subtract the two norm equations and put
\[
   u=(B'-A)/2,\quad v=(B'+A)/2,\quad
   N=q^6-p^6=uv.
\]
We have $0<u\le v$, $A=v-u$, and the positive integer
\[
                M=4p^6-(v-u)^2=x b^2
\]
is greater than $X_0$. Since $N=(q^3-p^3)(q^3+p^3)$ and $N/u-u$ decreases for $u>0$, positivity of $M$ is equivalent, for $u\le\sqrt N$, to $u>q^3-p^3$. Consequently it suffices to examine the positive divisors in
\[
       q^3-p^3<u\le\lfloor\sqrt{q^6-p^6}\rfloor .
\]
For an odd prime $r\le2000$, a nonzero square $-M\bmod r$ forces $r\nmid bD$ and $(D/r)=1$ because $-M=D b^2$. Such an $r$ would split in the candidate field. This contradicts $x>4\cdot2000^6$. It suffices to find one such witness for every remaining $M$; no extraction of its squarefree part is needed.

\paragraph{Finite lemma.} For each pair of primes $p<q\le5500$ satisfying $4p^6>X_0$, every divisor in the displayed interval yielding $M>X_0$ has a witness $r$ as above. The following counts include multiplicities over prime pairs and divisors.
\begin{center}
\begin{tabular}{lr}\toprule
Quantity & Count\\\midrule
Primes $p\le5500$ with $4p^6>X_0$ &414\\
Unordered pairs &85,491\\
Divisors in the strict-lower-endpoint interval &10,144,433\\
Pairs and divisors giving $M>X_0$ &9,393,302\\
Candidates with no witness &0\\\bottomrule
\end{tabular}
\end{center}
\textit{Verification.} Factor
\[
 N=(q-p)(q+p)(q^2+pq+p^2)(q^2-pq+p^2).
\]
Trial division is sufficient because each factor is at most $3\cdot5500^2$. Generate the divisors from their prime exponents, test the two inequalities and reduce $-M$ modulo the small primes. The Python implementation uses arbitrary-precision integers and sets of nonzero squares. The separately written C++ implementation uses 128-bit products, checks its factorization, and tests the equivalent residue $(v-u)^2-4p^6$ via precomputed masks. It includes the lower endpoint, which adds one divisor per pair with $M=0$ and affects none of the positive candidates. In this run primes from 19 to 151 already supplied a witness in every case. Both programs returned 9,393,302 tests above the threshold and zero survivors. The finite lemma contradicts the forced shared coefficient, proving the theorem.

\paragraph{The 432 fields and the order classification.}
The file \texttt{data/exponent-six.csv} lists every field discriminant and class number. For two through seven distinct ramified primes the counts in order are $60,141,140,76,12,3$. Aiken, Ross, Shvydkoy and Xue's Theorem~1.3(5) states that under ERH their table of exponent-six negative order discriminants is complete if EKN's remaining interval contains no fundamental example. Their construction reduces all analytic hypotheses to the underlying quadratic classifications. The integral endpoint $X_0$ is already covered by EKN (and is not a fundamental discriminant). Our theorem discharges the remaining condition. Thus exactly \textbf{1,236} imaginary quadratic orders have proper ideal class group of exponent six, of which $1,236-432=804$ are nonmaximal. The largest absolute order discriminant is $43\,522\,752$. We checked the 1,236 distinct entries in the authors' data and verified that its subset of 432 fundamental discriminants agrees entry for entry with the field table.
\href{https://github.com/erick-ross/research-data/blob/119301b8bcae42a891062d4877b77b6dbc8e820f/class-groups-small-exp/Table-6.txt}{Pinned order data}.
\href{https://arxiv.org/pdf/2605.26020}{Aiken, Ross, Shvydkoy and Xue, Theorem~1.3(5), Section~6 and Table~A.6}.
Using primitive positive reduced binary quadratic forms we verified the listed group orders and the sixth power of every one of the $103\,830$ classes in the 1,236 orders. More explicitly, after changing variables so that the leading coefficient $a$ is prime to the order discriminant, the form represents an invertible ideal $(a,\omega-u)$ in a quadratic order $\mathbb Z[\omega]$. Hensel-lifting the simple root $u$ of the polynomial of $\omega$ modulo $a$ gives $u_6$ modulo $a^6$. The ideal $(a^6,\omega-u_6)$ equals the sixth power by containment and equality of norms. Reduction of its norm form checks principality. Counting the reduced forms gives the order of the class group; a finite abelian group killed by six has elementary abelian Sylow subgroups, so its order determines its structure. This independently checks membership throughout the table. The cited conductor theorem supplies its exhaustiveness.

\subsection{What the same test gives at exponent seven}
For exact exponent seven, EKN's initial ceiling is $3.9\cdot10^{30}$. Outside their completed census the class number is odd, so genus theory gives one ramified prime. The only possibilities outside the small dyadic cases have $D=-x$ with $x$ prime. In particular there is no ramified prime in the test supports below. If no prime splits up to $Y$, the lower bound for $\log x$ becomes $S+2-14C/9$. We give the intermediate cutoffs to make the implication reproducible.
\begin{center}\begin{tabular}{rrr}\toprule
$Y$ & previous $\log x<$ & forced $\log x>$ if no prime splits\\\midrule
7000 & 70.50 & 76.0\\
5000 & 63.37 & 63.6\\
4700 & 61.01 & 61.5\\
4600 & 60.58 & 60.8\\\bottomrule
\end{tabular}\end{center}
Each contradiction produces a split $p\le Y$, so $x\le 4Y^7$. The final bound is $4\cdot4600^7=1.743270628864\cdot10^{26}$. EKN found 33 fields in the completed range; this calculation leaves the interval between their search and the new upper limit open.

\medskip\noindent\textbf{Exact arithmetic.}
If $v_i$ are the displayed heights, $F(j\delta)=
(\sum_i v_i v_{i+j})/(\sum_i v_i^2)$, taking heights to be zero
outside the twelve intervals. Interpolation between consecutive knots
is linear. The accompanying \texttt{zero\_test.py} compares
$(p^k)^{12}$ with powers of $Y$ to locate the interval exactly. It
encloses logarithms by reducing to $1\le z\le2$ and taking 20 terms
of the series for $2\operatorname{arctanh}((z-1)/(z+1))$, with its
geometric tail. Integer square roots give bounds on $1/\sqrt{p^k}$.
The code then verifies each finite prime sum, each slope bound and
the chain of contradictions using rationals. No zero ordinate or class
group is supplied as numerical input to these tests.

\par\medskip\section{Toroczkai's prime-gap graphs and simultaneous growth}
\label{sec:primegapgraphs}
Let $p_n$ be the $n$th prime, put $p_0=1$, and set $d_i=p_i-p_{i-1}$.
A prime-gap graph on $n\ge2$ vertices is a simple graph on $1,\ldots,n$
with $\deg(i)=d_i$. The sequence begins $1,1,2,2,4,2,4,2,4,6$.
Its degree sum is $p_n-1$. A matching consists of edges with pairwise
disjoint endpoints. Write $\nu(G)$ for its largest possible size.

\paragraph{Two named conjectures.}
Toroczkai conjectured that such a graph exists for every $n\ge2$, and
that \emph{every} such graph has a matching of size $d_{n+1}/2$.
Erd\H{o}s, Harcos, Kharel, Maga, Mezei and Toroczkai
proved both assertions under the classical Riemann hypothesis
(Theorems 1.1 and 1.6 of \cite{graphEH}).
The second implies the first. Begin with a single edge and replace
$d_{n+1}/2$ disjoint edges by paths of length two through a common new vertex.
The old degrees stay fixed and the new vertex has the required degree.
The zeta-function endpoint of the Lean repository therefore supplies the full
hypothesis of both published implications. The graph-theoretic results and
their analytic bridge are used here as theorems in print.
The RH-free bounds of Aggarwal, Frot, Gou and Wang
give existence for $n\ge\exp(\exp(30.32))$ and the matching assertion for
$n\ge\exp(\exp(34.33))$ \cite{graphAF}.

\subsection{A reserve of disjoint edges in every realization}
The same prime-gap estimate also allows many future insertions to be planned
using only the edges already present. All logarithms are natural.
Write $p=p_n$, $L=\log p$, and
\[
       N=\lceil652L^2\rceil,\qquad
       B(p)=\left\lceil\frac{p-2}{4N}\right\rceil .
\]
\paragraph{Matching and block theorem.}
Assume classical RH. Every prime-gap graph $G$ on $n\ge2$ vertices has
$\nu(G)\ge B(p)$. Whenever $p_{n+k}-p\le2B(p)$, the next $k$ vertices
can all be inserted using disjoint sets of old endpoints. No pair of new
vertices is adjacent, and every old vertex is incident to at most one
new edge. In particular, for each fixed $0<c<1/1304$ and all sufficiently
large $n$, this holds for $k=\lfloor c p/L^3\rfloor$.

\paragraph{Proof.}
For an integer $T\ge3$, let $H(T)=\sum_{i\le n,\ d_i\ge T}d_i$.
Theorem 1.8 of \cite{graphEH} bounds the sum of gaps of length at least $T$
starting at a prime in $[x,2x]$ by $163x(\log x)^2/T$ under RH,
for every real $x\ge2$. Apply this at $x=p/2^j$ for $j\ge1$ while $x\ge2$.
These intervals cover every gap start up to $p$ except possibly $1,2,3$,
whose following gaps have size at most two. Since $\sum_{j\ge1}p/2^j=p$,
\[
                  H(T)\le 163pL^2/T.                 \tag{G1}\label{eq:g1}
\]
Delete all vertices of degree at least $T$. This loses at most $H(T)$
edges, leaving maximum degree at most $T-1$. Vizing's edge-coloring
theorem partitions the remaining edges into at most $T$ matchings.
In particular
\[
 \nu(G)\ \ge\
 \left\lceil \frac{(p-1)/2-H(T)}{T}\right\rceil .    \tag{G2}\label{eq:g2}
\]
Taking $T=N$ gives $H(N)\le p/4$, hence the advertised $B(p)$.
The constant comes from the leading lower bound
$p/(2T)-163pL^2/T^2$, which is maximal at $T=652L^2$.

Partition a matching of size $(p_{n+k}-p)/2$ into pieces of respective sizes
$d_{n+1}/2,\ldots,d_{n+k}/2$. For each piece, remove its edges and join
their endpoints to the corresponding new vertex. Distinct pieces have
disjoint endpoints. This implements the claimed insertions in any order.
Finally $2B(p)\sim p/(1304L^2)$. For fixed $c<c'<1/1304$ take
$h=c'p/L^2$. Schoenfeld's RH estimate \cite[Corollary 1]{graphS}
$\pi(t)=\operatorname{li}(t)+O(\sqrt t\log t)$ gives
$\pi(p+h)-\pi(p)=(c'+o(1))p/L^3$ by subtraction at the two endpoints.
Thus the first $\lfloor c p/L^3\rfloor$ new primes lie within the allowed
distance once $p$ is sufficiently large. \hfill$\square$

\paragraph{The multiset extension question.}
Section 6 of \cite{graphDPG} asks which multisets of prescribed degrees
can be appended to a fixed initial graph by degree-preserving steps.
For a list of positive even degrees $e_1,\ldots,e_k$, the restricted
version using only original edges with pairwise disjoint endpoints has
the exact criterion $\sum_j e_j\le2\nu(G)$. Necessity follows because
all removed edges form one matching. Sufficiency is the partition
construction above. The unrestricted problem permits reusing vertices
and edges created by earlier steps. For example, starting with a single
edge one can successively add arbitrarily many vertices of degree two
by subdividing one current edge at each step. The block theorem gives
a quantitative sufficient condition for the unrestricted problem on
every prime-gap realization, together with the stronger guarantee on
which edges the extensions use.

\paragraph{Exact example.}
For $n=10{,}000$, $p=104729$ and $T=31$, the actual degree sum above
the cutoff, counting equality, is $9226$. Equation (G2) grants
$\lceil(52364-9226)/31\rceil=1392$ disjoint edges in \emph{every}
realization, without an analytic hypothesis. The next 236 primes end
at $107509$ and require $(107509-104729)/2=1390$ edges in total.
The following prime is $107563$. The accompanying Python code checks the
degree sums by sieving, constructs one realization by Havel--Hakimi,
and carries out all 236 insertions. It verifies simplicity, all $10236$
final degrees, and distinctness of the old endpoints.

\subsection{Keeping every late edge close in vertex labels}
An infinite sequence of degree-preserving steps can keep moving the same
edge endpoint indefinitely. A choice that stops moving old endpoints gives
a connected countable graph with the prescribed degrees. The following
result also controls how long each vertex has to wait for its final neighbors.

\paragraph{Stabilization theorem.}
Assume classical RH. There is a sequence of connected prime-gap graphs
$G_2,G_3,\ldots$, obtained by the insertion operation above, in which
the neighbors of every fixed vertex eventually become permanent.
Their limit is a connected, locally finite simple graph on the positive
integers with $\deg(i)=d_i$. Every limit edge $ij$ with $i<j$ satisfies
$j-i=O(j/\log j)$ as $j\to\infty$, uniformly in $i$.

\paragraph{Proof.}
Fix $A=100$ and a sufficiently large constant $P\ge e^{A}$. Use the
published all-$n$ theorem until the first stage with $p=p_n\ge P$.
Connectivity persists because each replaced edge has a new path
between its original endpoints. At that stage choose a spanning tree.
Thereafter all of its edges, and the edges of the successive trees
defined below, will be protected. Set
\[
 w(x)=\begin{cases}
 x,&0\le x<P,\\
 Ax/(A+\log(x/P)),&x\ge P,
 \end{cases}
 \qquad F(x)=x-w(x).
\]
At stage $n$ require that no edge used in the next insertion touch
any of the first $k=\pi(F(p))$ vertices. Both $w$ and $F$ are
nondecreasing, and $F$ is 1-Lipschitz. For $p\ge P$,
$w(p)\ge Ap/\log p$.

We show inductively that enough eligible edges remain. Let $u=p_k$ if
$k>0$. At stage $k$, all neighbors of the first $\pi(F(u))$ vertices
had labels at most $k$ and were permanently protected from that stage
onward by monotonicity of the cutoff. If $u<P$, this assertion is empty.
Thus at least $(v-1)/2$ edges inside the first $k$ labels still exist
at stage $n$, where $v$ is the greatest prime at most $F(u)$, using $v=1$
if there is none. This follows by summing the degrees of those earlier
vertices and allowing each internal edge to be counted twice.
The degree sum in the entire first $k$ labels is $u-1$, so the number of
edges incident to them is at most $(u-1)-(v-1)/2$.
Consequently at least $(p-2u+v)/2$ edges have both endpoints outside
the first $k$ labels.

By the RH prime-gap bound $d_{m+1}=O(\sqrt{p_m}\log p_m)$
(Theorem 2.3 of \cite{graphEH}), $v\ge F(u)-O(\sqrt p\log p)$.
As $u\le F(p)$ and $F$ is 1-Lipschitz,
$-2u+F(u)\ge-2F(p)+F(F(p))$. Hence the number of available edges
before protecting the tree or truncating high degrees is at least
\[
 \frac{p-2F(p)+F(F(p))}{2}-O(\sqrt p\log p)
 =\frac{w(p)-w(F(p))}{2}-O(\sqrt p\log p).
\]
For $k=0$ the same inequality follows directly from the total
$(p-1)/2$ edges and $0\le F(p)<2$.
For $P\le x\le p$, differentiation gives
$w'(x)=A(t-1)/t^2$ with $t=A+\log(x/P)\ge A$, hence
$w'(x)\ge w(p)/(2p)$. For $x<P$ the derivative equals one.
Integrating on $[F(p),p]$ gives
\[
                      w(p)-w(F(p))\ge \frac{w(p)^2}{2p}.
\]

The part of the protected tree induced on the $n-k$ remaining vertices
is a forest and contains at most $n-k$ edges. For all sufficiently large
$p$, uniformly in $P\le p$, the RH prime-counting estimate gives
$n-k\le3w(p)/\log p$. Indeed for $w(p)\ge p/2$ use
$\pi(p)\sim p/\log p$. Otherwise subtract the estimates for $\pi$ at
$p$ and $p-w(p)\ge p/2$ and use $w(p)\ge Ap/\log p$ to absorb the error.
Also exclude vertices of degree at least $T=\lceil(\log p)^5\rceil$.
By (G1) the resulting loss is at most $163p/(\log p)^3$ edges.
The number left is thus at least
\[
 \frac{w(p)^2}{4p}-\frac{3w(p)}{\log p}
       -O(\sqrt p\log p)-\frac{163p}{(\log p)^3}
          \gg \frac{p}{(\log p)^2}.
\]
Here $w(p)^2/(4p)\ge25w(p)/\log p$. All implied constants are
independent of $P$, so a single sufficiently large $P$ makes these
estimates valid throughout the induction. Edge coloring now gives
a matching of size $\gg p/(\log p)^7$, greater than $d_{n+1}/2$.
Make the insertion and add one edge incident to the newcomer to the
previously protected tree. The resulting trees are nested.

Since $F(p)\to\infty$, the neighbors of each fixed label eventually
become permanent. The union of the trees proves connectivity of the
limit, and stabilization preserves every finite degree. An edge $ij$,
$i<j$, is born at stage $j$, so for $p=p_{j-1}$ we have
$i>\pi(F(p))$ once $p\ge P$. With $P$ fixed,
$w(p)\sim100p/\log p$ and
$j-i\le1+\pi(p)-\pi(F(p))=O(p/(\log p)^2)=O(j/\log j)$.
\hfill$\square$

\paragraph{Which part uses RH?}
A sufficiently late starting stage for the stabilization theorem can
also be treated with older unconditional estimates. Theorems 2.2 and
2.4 of \cite{graphEH} give $d_{n+1}\le p^{5/8+o(1)}$ and
$\sum_{i\le n}d_i^2\le p^{4/3+o(1)}$. Use
$T=\lceil p^{7/20}\rceil$, so
$H(T)\le T^{-1}\sum d_i^2\le p^{59/60+o(1)}=o(p/(\log p)^2)$.
The gap estimate makes the error from rounding $F(u)$ down to a prime
$o(p/(\log p)^2)$ as well. The classical prime number theorem with
remainder $O(x/(\log x)^4)$ suffices for $n-k\le3w(p)/\log p$.
After the same exclusions, a matching has size
$\gg p^{13/20}/(\log p)^2$, exceeding $d_{n+1}/2$.
Thus any sufficiently late connected starting realization admits this
construction using classical bounds alone. The all-stage start at two
vertices comes from the published RH implication. The uniform matching
bound with $T=\lceil652L^2\rceil$ uses the RH tail estimate (G1).

\par\medskip\section{The discriminant and conductor of the reflected fields}
\label{sec:local}
The first incoming report ends its degree-five classification by asking for the ramification at five. The two possible exponents already appear in the reflection theorem of Cohen, Rubinstein-Salzedo and Thorne. To decide which exponent occurs for a given CM field, we compute the local power class of the generator used in that report. This also gives the conductor in the Galois closure and distinguishes the fields over the two quadratic bases with class number 25.
\href{https://doi.org/10.1112/S0010437X15007411}{Cohen, Rubinstein-Salzedo and Thorne (2015)}.
\href{https://thornef.github.io/countDLgen-final.pdf}{Cohen and Thorne, Theorem~9.7}.

\subsection*{Which fields the computation exhausts}

Let $D<0$ be a fundamental discriminant, let $\ell$ be an odd prime
coprime to $D$, and put $E=\mathbb Q(\sqrt D)$ and $B=E(\zeta_\ell)$.
Write $r=\dim_{\mathbb F_\ell}\operatorname{Cl}(E)[\ell]$.
The number of fields in the reflection family attached to $D$ is
$(\ell^r-1)/(\ell-1)$. This counts isomorphism classes, and its
meaning includes a specified Galois action.

To see the action, take a nonzero class $[\mathfrak a]$ of order
$\ell$, write $\mathfrak a^\ell=(\alpha)$ and put
$n=N\mathfrak a>0$. Then $\alpha\bar\alpha=n^\ell$.
Changing the ideal representative changes $\alpha$ by an $\ell$th
power, up to a unit. Every unit of $E$ is an $\ell$th power under the
stated hypotheses. Thus the virtual units
\[
 \{v\in E^\times:(v)\ \text{is an }\ell\text{th ideal power}\}/E^{\times\ell}
 \simeq \operatorname{Cl}(E)[\ell]
\]
identify ideal-class lines with radical lines. The radical line remains
nonzero and distinct from the other lines in $B^\times/B^{\times\ell}$
because restriction followed by norm raises classes to $[B:E]=\ell-1$.

Set $M=B(\beta)$ with $\beta^\ell=\alpha$. A generator $s$ of
$\operatorname{Gal}(M/B)$ sends $\beta$ to $\zeta_\ell\beta$.
For a primitive root $g$ modulo $\ell$ choose $u$ fixing $E,\beta$
and sending $\zeta_\ell$ to $\zeta_\ell^g$. Complex conjugation $c$
sends $\beta$ to $n/\beta$ and $\zeta_\ell$ to $\zeta_\ell^{-1}$.
These automorphisms generate the full Galois group. Direct substitution
shows that $c$ is central and $usu^{-1}=s^g$.
The $\ell$ real conjugates of $\beta+n/\beta$ are
$\zeta_\ell^j\beta+n/(\zeta_\ell^j\beta)$.
They are distinct. Equality for different $j$ would give
$\beta^2/n$ an $\ell$th root of unity and hence $\alpha^2=n^\ell$,
forcing $[\alpha]=0$ in $E^\times/E^{\times\ell}$.
Their splitting field is $M^{\langle c\rangle}$,
since the affine action of $s$ and $u$ on their indices is faithful.
Its cyclic degree-$(\ell-1)$ subfield $B^{\langle c\rangle}$ is the
mirror field in the sense of Cohen, Rubinstein-Salzedo and Thorne.
The displayed conjugation rule is exactly their star condition.

Two different ideal-class lines give different normal closures.
On adjoining $E$ each normal closure recovers its Kummer extension
$M/B$, and Kummer theory distinguishes the lines. Different generators
of one line give the same normal closure. Its cyclic complements of
order $\ell-1$ are conjugate in the affine group, so they give a single
isomorphism class of degree-$\ell$ fields.

For each of the discriminants used here, the local discriminant
calculation places these fields among the totally real starred fields
of discriminant $|D|^{(\ell-1)/2}\ell^{\ell-2}$ or
$|D|^{(\ell-1)/2}\ell^\ell$.
Theorem~1.2 and Lemma~2.8 of Cohen, Rubinstein-Salzedo and Thorne
give precisely $(\ell^r-1)/(\ell-1)$ starred fields in the union
of these two sets. The construction therefore exhausts this family.
The discriminant alternatives and the pooled count are established
in that paper. The present table assigns a specific alternative and
defining polynomial to each ideal-class line.
\href{https://arxiv.org/pdf/1405.1075v2}{Cohen, Rubinstein-Salzedo
and Thorne, Theorem~1.2, Lemma~2.8 and Proposition~3.10}.

The analytic input enters in exhausting the base fields.
Elsenhans, Kl\"uners and Nicolae, Theorem~1 and Table~1, find 27
imaginary quadratic fields of class-group exponent five through
$3.1\cdot 10^{20}$. Their Lemma~1 and Theorem~3 bound every possible
absolute discriminant of exponent five by $2.3\cdot 10^{20}$ under the
quadratic least-split-prime estimate (Table~2). The requisite critical-line
statement for $\zeta_E(s)=\zeta(s)L(s,\chi_D)$ follows from the public
Dirichlet GRH theorem. Their published data give 25 groups $C_5$ and
two groups $C_5^2$, at $D=-12451,-37363$. Thus the count is
$25+2(5^2-1)/(5-1)=37$.
\href{https://arxiv.org/pdf/1803.02056}{Elsenhans, Kl\"uners and
Nicolae, Theorems~1--3 and Tables~1--2}.
\href{https://math.uni-paderborn.de/en/ag/ca/research/exponent/exponent0}
{Authors' tables of discriminants and groups}.


\subsection{A local test with two possible answers}
For the calculation assume $D=-q$ is an odd prime discriminant and $q\ne\ell$. Use the notation $E,B,n,\alpha,\beta$ above and set
\[
 K=\mathbb Q(\beta+n/\beta),\qquad C=M^{\langle c\rangle},\qquad F=B^{\langle c\rangle}.
\]
Here $C$ is the normal closure of $K$ and $C/F$ is cyclic of degree $\ell$. For clarity, the conductor below is the relative Artin conductor of any nontrivial character of $C/F$. Let $v\mid\ell$ in $E$, $U=E_v$, let $f=[U:\mathbb Q_\ell]$, and let $u$ be the unit obtained from $\alpha$ by dividing by $\ell^{v(\alpha)}$, with $v(\ell)=1$. Since $\ell\mid v(\alpha)$, this operation preserves the local radical class. Then
\[
 \begin{array}{c|c|c}
  \text{local test in }\mathcal O_U/\ell^2&
  |\operatorname{disc}(K)| & \text{conductor of } C/F\\ \hline
  u^{\ell^f-1}=1 & q^{(\ell-1)/2}\ell^{\ell-2} & (1) \\
  u^{\ell^f-1}\ne1 & q^{(\ell-1)/2}\ell^{\ell} & \lambda^2
 \end{array}
\]
Here $\lambda$ is the unique prime over $\ell$ in $F$. If $\ell$ splits in $E$, the two choices of $v$ give the same answer since $\alpha\bar\alpha=n^\ell$.

\paragraph{The local algebra at $\ell$.}
Since $K E=E(\beta)$, extension of scalars gives
\[
 (K\otimes_{\mathbb Q}\mathbb Q_\ell)\otimes_{\mathbb Q_\ell}U
       \simeq U[X]/(X^\ell-u).
\]
Here $KE=E(\beta)$ follows also from $[E(\beta):E]=\ell$ and $\beta^2-(\beta+n/\beta)\beta+n=0$. The unramified extension $U/\mathbb Q_\ell$ commutes with formation of the maximal order of an \'etale algebra: on each local factor its base change is a product of rings of integers of unramified extensions. Hence the discriminant valuation can be read after this scalar extension.

Decompose a unit into a Teichm\"uller root of unity of order dividing $\ell^f-1$ and a principal unit. The $\ell$th power map is an automorphism on the former group. On the latter the logarithm identifies $1+\ell\mathcal O_U$ with $\ell\mathcal O_U$, and taking an $\ell$th power has image $1+\ell^2\mathcal O_U$. This proves the displayed residue test. In the power case the algebra is $U\times U(\zeta_\ell)$ with a tame, totally ramified factor of degree $\ell-1$. Its discriminant exponent is $\ell-2$.

In the other case take $z$ to be any lift of the unique residue root of $X^\ell-\bar u$. Then $v(z^\ell-u)=1$, since a solution modulo $\ell^2$ would pass the power test. The shifted polynomial $(X+z)^\ell-u$ is Eisenstein. Its root is a uniformizer and generates the integer ring; evaluating the derivative at the root of the original polynomial gives $\ell$ times a unit. Its norm has valuation $\ell$. This proves the wild exponent.

\paragraph{The remaining primes and the conductor.}
For a finite prime of $B$ away from $\ell$, the valuation of $\alpha$ is a multiple of $\ell$. After removing a local $\ell$th power one adjoins a root of a unit polynomial $X^\ell-v_0$ with unit discriminant, hence an unramified extension. Consequently ramification away from $\ell$ can occur only at $q$. There inertia has order two in $M/\mathbb Q$, fixes $\zeta_\ell$, and acts on the $\ell$ real roots by $j\mapsto h-j$ for some $h$. It has one fixed point and $(\ell-1)/2$ transpositions, so tame inertia contributes $(\ell-1)/2$ to the discriminant exponent in $K$. This inertia also injects into $\operatorname{Gal}(F/\mathbb Q)$, so $C/F$ is unramified there.

At $\ell$, inertia acts on $F$ through the full cyclotomic group, since $E$ is unramified there and the cyclotomic group acts faithfully on $F$. Thus $F/\mathbb Q$ is totally tamely ramified at $\ell$. In the power case adjoining $U$ to the local splitting field leaves it inside $U(\zeta_\ell)$, of absolute degree at most $2(\ell-1)$. Its local degree over $F_\lambda$ divides $\ell$, so $\lambda$ splits completely in $C/F$.

In the wild case, over $U$ the splitting field contains the degree-$\ell$ totally ramified radical extension and the degree-$(\ell-1)$ totally ramified cyclotomic extension. Indeed the normal closure of $K$ after this base change is the splitting field of the isomorphic radical algebra. The ramification index is divisible by both degrees. Its upper bound $\ell(\ell-1)$ shows that the entire affine group is inertia already over $\mathbb Q_\ell$. Use $K$ locally in the discriminant tower. The extension $C/K$ is tame of degree $\ell-1$, so
\[
 v_\ell(d_{C/\mathbb Q})=(\ell-1)\ell+(\ell-2)=\ell^2-2.
\]
Using $F$ instead gives $v_\lambda(d_{C/F})=\ell^2-2-\ell(\ell-2)=2(\ell-1)$, as all residue degrees are one. The conductor--discriminant formula for the cyclic group of prime order gives conductor exponent two for each of its $\ell-1$ nontrivial characters. This proves the conductor assertion.

\subsection{All 26 coefficient fields from the incoming quintic classification}
For the following 25 rows the given $p$ is a split rational prime, $4p^5=a^2+qb^2$, and $p\nmid a$. Set $\alpha=(a+b\sqrt{-q})/2$. Its ideal is the fifth power of one of the primes above $p$. As $p<q/4$, a generator of norm $p$ cannot exist. The associated field has defining polynomial
\[
 P_{p,a}(X)= X^5-5pX^3+5p^2X-a.
\]
This follows by expanding $(\beta+p/\beta)^5$ and using $\beta^5+\bar\beta^5=a$. The table gives $v_5=v_5(\operatorname{disc}K)$; the full discriminant is $q^2 5^{v_5}$. Each half of the table reads downward.
\begin{center}\small
\begin{tabular}{rrrrr@{\qquad}rrrrr}
\toprule
$q$ & $p$ & $a$ & $b$ & $v_5$ & $q$ & $p$ & $a$ & $b$ & $v_5$\\
\midrule
47 & 2 & 9 & 1 & 5 & 691 & 5 & 38 & 4 & 5\\
79 & 2 & 7 & 1 & 5 & 739 & 5 & 26 & 4 & 3\\
103 & 2 & 5 & 1 & 5 & 787 & 7 & 59 & 9 & 5\\
127 & 2 & 1 & 1 & 5 & 947 & 3 & 5 & 1 & 5\\
131 & 3 & 29 & 1 & 5 & 1051 & 5 & 107 & 1 & 3\\
179 & 3 & 16 & 2 & 5 & 1123 & 7 & 239 & 3 & 3\\
227 & 3 & 8 & 2 & 5 & 1723 & 11 & 149 & 19 & 3\\
347 & 3 & 25 & 1 & 3 & 1747 & 17 & 1565 & 43 & 5\\
443 & 3 & 23 & 1 & 3 & 1867 & 7 & 4 & 6 & 5\\
523 & 7 & 220 & 6 & 5 & 2203 & 7 & 255 & 1 & 5\\
571 & 5 & 58 & 4 & 5 & 2347 & 17 & 1810 & 32 & 5\\
619 & 5 & 109 & 1 & 5 & 2683 & 11 & 344 & 14 & 5\\
683 & 3 & 17 & 1 & 5 &  & & & & \\
\bottomrule\end{tabular}
\end{center}

\noindent The remaining CM coefficient field is $\mathbb Q(\zeta_{11})^+$, with discriminant $11^4$, by the conductor--discriminant formula for its four nontrivial characters. Thus six of the 25 noncyclotomic fields are tame at five. The absolute CM discriminants of these six are $347,443,739,1051,1123,1723$.

\paragraph{Worked computation.}
For $q=347$ one may take $p=3$, $a=25$, $b=1$. Modulo 25 we have $(2\sqrt{-347})^5=(25+\sqrt{-347})/2$, since $32\cdot347^2\equiv13\pmod{25}$. This gives field discriminant $5^3\cdot347^2=15\,051\,125$ and index five. For $q=47$, $p=2$, $a=9$, $b=1$, the field discriminant is $5^5\cdot47^2=6\,903\,125$ and the displayed polynomial generates the full integer ring. In fact the polynomial discriminant is $5^5(4p^5-a^2)^2=5^5q^2 b^4$, by a resultant with $P'_{p,a}$ (or the critical values $\pm2p^{5/2}$ of the Dickson polynomial). Thus the power-basis index is $5b^2$ in the tame case and $b^2$ in the wild case.

\subsection{Twelve fields from two rank-two class groups}
Over each of $D=-12451$ and $D=-37363$ the six lines of $\mathbb F_5^2$ give six nonisomorphic reflected quintics. Exactly one line in each group is tame. These twelve fields extend the reflection computation; the weight-two CM classification in the incoming report uses the preceding 25 rows and the cyclotomic exception. All five wild fields over any one rank-two base have exactly the same discriminant.
\begin{center}\small
\begin{tabular}{rrrrr@{\qquad}rrrrr}
\toprule
\multicolumn{5}{c}{$D=-12451$} & \multicolumn{5}{c}{$D=-37363$}\\
$p$ & $a$ & $b$ & label & $v_5$ & $p$ & $a$ & $b$ & label & $v_5$\\
\midrule
5 & 7 & 1 & $[2\!:\!0]$ & 3 & 11 & 779 & 1 & $[3\!:\!1]$ & 5\\
7 & 132 & 2 & $[1\!:\!1]$ & 5 & 13 & 942 & 4 & $[1\!:\!3]$ & 5\\
11 & 577 & 5 & $[2\!:\!3]$ & 5 & 23 & 4607 & 11 & $[3\!:\!3]$ & 5\\
13 & 1134 & 4 & $[1\!:\!3]$ & 5 & 37 & 13131 & 53 & $[4\!:\!1]$ & 5\\
53 & 23339 & 301 & $[0\!:\!1]$ & 5 & 43 & 7677 & 119 & $[0\!:\!1]$ & 3\\
109 & 79839 & 2105 & $[2\!:\!4]$ & 5 & 79 & 96628 & 282 & $[3\!:\!0]$ & 5\\
\bottomrule\end{tabular}
\end{center}


Here a projective label certifies that its row represents a distinct line. For each base choose two triples $(s,r,z)$ where $s\equiv1\pmod5$ is prime, $r^2\equiv D\pmod s$ and $z$ has order five. Reducing $\alpha$ by $\sqrt D\mapsto r$ and raising to $(s-1)/5$ gives $z^e$, with $e\in\mathbb F_5$. The two exponents form the label $[e_1:e_2]$. All table generators are units at the probe primes. A relation between their classes modulo fifth powers would remain a relation between the labels because the residual fifth-power characters are multiplicative. The six nonproportional labels prove exhaustion by the known rank. The triples are $(181,20,59),(211,93,107)$ for the first base and $(151,38,59),(241,76,87)$ for the second.

The checker enumerates all 625 potential roots modulo 25 in the quadratic ring. When $p=5$ the quadratic algebra at five splits; one component of $\alpha$ is a unit, the other has valuation five. Modulo 25 the latter component is zero. The unit component decides the test, and its reciprocal times $5^5$ gives the other local class. A separate computation of maximal orders using SymPy's Round Two algorithm verifies every polynomial's field discriminant and order index, including all twelve rows above.

\par\medskip\section{One prime-splitting test for six Dedekind zeta functions}\label{sec:joint-test}

Fix $D=-12451$ or $-37363$. Let $K_1,\ldots,K_6$ be the six reflected
quintic fields for this $D$. At an unramified prime $p$, choose any two
distinct fields $K_i,K_j$ and record whether $p$ splits completely in
both. The same primes pass this test for every choice of the pair.
Together with a residue-class condition, this single bit gives a
one-sided criterion for RH in all six fields.

\paragraph{Joint Galois group.}
Choose virtual units $\alpha_1,\alpha_2$ representing a basis of
$\operatorname{Cl}(E)[5]$. Over $B=E(\zeta_5)$ their radical classes
are independent by restriction and norm. Adjoining both fifth roots
gives $M/B$ with translation group $V=\mathbb F_5^2$ by Kummer theory.
Choose the cyclotomic automorphisms fixing $E$ and both radicals,
and let $c$ be complex conjugation. As in the single-line case $c$ is
central. The six line evaluations of $V$ separate points, and each
quintic splitting field detects the full scalar group. Thus the joint
normal closure $T$ is $M^{\langle c\rangle}$ and
\[
 \operatorname{Gal}(T/\mathbb Q)
   = V\rtimes\mathbb F_5^\times,\qquad
 H=T^V=B^{\langle c\rangle}.
\]
The scalar group multiplies both coordinates. On the five embeddings
of the field associated with a radical line, a group element $(v,a)$
acts by $j\mapsto aj+\lambda(v)$, where $\lambda$ is a representative of the corresponding
line in the dual space. Its scalar at $p\nmid5D$ is
$\mu(p)=\chi_D(p)(p\bmod 5)\in\mathbb F_5^\times$.
Indeed Frobenius acts on $\zeta_5$ by raising to $p$ and on $E$ via
$\chi_D$. Simultaneous complex conjugation induces the identity on $H$.

For $a\ne1$ the affine map is conjugate to multiplication by $a$.
For $a=1$ a nonzero translation acts as a five-cycle. Of the six
functional lines, exactly one vanishes on a specified $v\ne0$.
Chebotarev gives the following natural densities; the second column
describes residue degrees.
\[
\begin{array}{lll}
 (v,a) & \text{splitting across the six fields} & \text{density}\\
 \hline
 v=0,\ a=1 & (1,1,1,1,1)\text{ in all six} & 1/100\\
 v\ne0,\ a=1 & (1,1,1,1,1)\text{ in one, }(5)\text{ in five}
      & 24/100\\
 a=-1 & (1,2,2)\text{ in all six} & 25/100\\
 a=2\text{ or }3 & (1,4)\text{ in all six} & 50/100
\end{array}
\]
Each of the six choices of the uniquely splitting quintic in the
second row occurs with density $4/100$. These densities hold
unconditionally for the two specified $D$.

\paragraph{The test.}
Let $S$ be a fixed finite set of excluded primes, containing those
dividing $5D$. When using reductions of defining polynomials one may
also exclude primes dividing their polynomial discriminants. Put
\[
 W(p)=25\,\mathbf1_{\{p\text{ splits completely in }K_i,K_j\}}
               -\mathbf1_{\{\mu(p)=1\}}, \qquad
 S_W(x)=\sum_{\substack{p\le x\\p\notin S}}
           W(p)\log p\,(p^{-1}-x^{-1}).
\]
The values of $W$ in the four rows are $24,-1,0,0$. There exists a
constant $c_W=\lim_{x\to\infty}S_W(x)$. Given Dirichlet GRH, the
following statements are equivalent.
\begin{enumerate}
\item All six $\zeta_{K_i}$ satisfy RH.
\item $\zeta_T$ satisfies RH.
\item For every $\epsilon>0$ there exists $C_\epsilon>0$ such that
$S_W(x)-c_W\le C_\epsilon x^{-1/2+\epsilon}$ for all sufficiently large $x$.
\end{enumerate}
Under these conditions, $|S_W(x)-c_W|\ll_{D,S} x^{-1/2}\log^2(2x)$.

\paragraph{Factorization and forward bound.}
Let $\tau_i$ be the character of the permutation representation on
the five embeddings of $K_i$ with one trivial character removed.
The table gives $W=\sum_i\tau_i$ on all Frobenius classes.
The regular character of $\operatorname{Gal}(T/\mathbb Q)$ equals the
inflated regular character of $\operatorname{Gal}(H/\mathbb Q)$ plus
$4\sum_i\tau_i$. Indeed their sum has value $100$ at the identity and
zero elsewhere. Artin formalism for permutation characters gives
\[
 \zeta_T(s)=\zeta_H(s)F(s)^4,\qquad
 F(s)=\prod_{i=1}^{6}\frac{\zeta_{K_i}(s)}{\zeta(s)}.
\]
Since $H/\mathbb Q$ is abelian, its zeta function factors into Dirichlet
$L$-functions. Dirichlet GRH makes each ratio in $F$ holomorphic on
$\Re s>1/2$, with a nonzero removable value at one, and gives the needed
zero-free region for $\zeta_H$ away from its pole. Zeros of the six
numerators in the right half of the strip therefore contribute positive
multiplicities to the product. The Dedekind functional equations give
the equivalence of the first two assertions.

Write $\theta_W(t)=\sum_{p\le t,\ p\notin S}W(p)\log p$. For each
$K_i$ the permutation trace counts degree-one prime ideals above $p$.
The fixed-field prime ideal theorem, applied to these six fields and
$\mathbb Q$, yields $\theta_W(t)=O(t\exp(-b\sqrt{\log t}))$ for some
$b>0$. The six main terms cancel. Higher residue degrees and prime
powers contribute $O_{D,S}(\sqrt t\log^2(2t))$, by summing over
$p^j\le t$, $j\ge2$. A possible real exceptional zero of a fixed field
is absorbed in the displayed asymptotic error. Under the first assertion,
the resulting bound is $\theta_W(t)=O_{D,S}(\sqrt t\log^2(2t))$.
Termwise integration gives
\[
 c_W=\int_1^\infty \theta_W(t)t^{-2}\,dt,\qquad
 R(x):=S_W(x)-c_W=-\int_x^\infty\theta_W(t)t^{-2}\,dt.
\]
This proves existence and the asserted two-sided estimate.
\href{https://arxiv.org/abs/1312.4463}{Greni\'e and Molteni,
\emph{Explicit versions of the prime ideal theorem for Dedekind zeta functions under GRH}}.

\paragraph{Converse from an upper bound.}
For $\Re s>1$ let $A(s)=\sum_{p\notin S}W(p)(\log p)p^{-s}$.
Euler factors give $A=-F'/F-J$ with $J$ holomorphic on $\Re s>1/2$.
At good primes the terms with powers $p^j$, $j\ge2$, are bounded by
$24(\log p)p^{-j\Re s}$ and converge locally uniformly.
The local factors at the finitely many excluded primes have inverse
roots of absolute value at most one, so their corrections are
holomorphic on $\Re s>0$. The prime ideal theorem gives $A(1)=c_W$
by partial summation. The Mellin transform of the centered sum is
\[
 P(s)=\int_1^\infty R(x)x^{-s}\,dx
     =\frac{A(s)}{s(s-1)}-\frac{c_W}{s-1}\qquad(\Re s>1).
\]
It is regular at one. For $\Re s>1/2$, any zero $\rho$ of $F$ of order
$m$ gives a simple pole of $P$ with residue $m/(\rho(1-\rho))$.
All its real poles in this region have positive residue, since a real
zero there lies below one.

Assume the third assertion and fix $1/2<\gamma<1$. Choose $C,X$
so $U(x)=Cx^{\gamma-1}-R(x)\ge0$ for $x\ge X$.
The Mellin integral $\int_X^\infty U(x)x^{-s}\,dx$ agrees with
$C/(s-\gamma)-P(s)$ up to an entire function, and its abscissa of
convergence is at most one. If that abscissa exceeded $\gamma$,
Landau's lemma for a nonnegative Mellin integrand would force a
singularity at the real abscissa. The meromorphic continuation permits
only a simple pole of negative residue there. Approaching the pole
from the right would make the integral negative. This contradiction
puts the abscissa at most at $\gamma$, so the integral is holomorphic
throughout $\Re s>\gamma$, where $P$ has no poles. Taking all
$\gamma>1/2$ proves zero-freeness of $F$ there. This proves the
converse. The prime estimate itself remains to be proved for these
nonabelian fields.

\paragraph{The information missing from residue classes.}
For every positive integer $m$ and every complex-valued function $h$
on $(\mathbb Z/m\mathbb Z)^\times$, Chebotarev gives
\[
 \lim_{x\to\infty}\frac1{\pi(x)}
 \sum_{\substack{p\le x\\p\notin S,\ p\nmid m}}
       |W(p)-h(p\bmod m)|^2
  =6+\frac1{\varphi(m)}
           \sum_{a\in(\mathbb Z/m\mathbb Z)^\times}|h(a)|^2 .
\]
Indeed $[V\rtimes\mathbb F_5^\times,V\rtimes\mathbb F_5^\times]=V$,
since the commutators with scalar $2$ give every translation. Thus
$T\cap\mathbb Q(\zeta_m)$ is contained in $H$. Conditional on a
cyclotomic Frobenius, the compatible elements of $\operatorname{Gal}(T/\mathbb Q)$
are uniformly counted on a union of full translation cosets. The mean
of $W$ on each coset is zero: on $V$ its sum is $24-24=0$, and on
the other three cosets $W$ vanishes. The mixed term in the square
therefore has zero mean. The first two rows of the table give the
unconditional mean square $(24^2+24)/100=6$.

The cubic case appears in
\href{https://github.com/VladimirReshetnikov/ProveIt/tree/1e1728ad221d69e48fea09bf754b73422af44943/docs/incoming}
{the incoming cubic-field report}. The Mellin argument above uses
\href{https://github.com/lsernau/grh-lean/blob/72fead102f663d5c7a407964a13dc0cf9d977207/grh-compact/src/New/Endgame/OneSidedCriterion.lean}
{the Lean repository's Landau criterion}.

\paragraph{Finite splitting check.}
The checker counts the roots of all six quintics over each good prime field through 5000, omitting $2,5$ and the prime divisors of the quadratic and polynomial discriminants. For $D=-12451$ it checks 663 primes. The first with all six splitting is 269; at 11 only the quintic with $p=7$ in the table splits completely. For $D=-37363$ it checks 661 primes and the first common splitting prime is 769. The first representative of every splitting case is also factored independently. The algebraic argument gives the densities, and the finite data illustrate the distinctions on which the analytic test depends.

\section{Further research questions}
\begin{enumerate}
\item Can the exponent-seven search be closed below $4\cdot4600^7$? For seventh powers, $q^7-p^7$ has a different factorization and a larger bound for the repeated coefficient. A usable extension would need to prove that two small split primes yield compatible norm equations or find a replacement elimination. This is the one exponent among $1,\ldots,8$ whose global census remains unresolved here.
\item For prime-gap graphs, what is the asymptotic minimum of the largest matching among all realizations on $n$ vertices? The block theorem gives a lower bound of order $p_n/(\log p_n)^2$; a better bound lengthens the block of insertions that use disjoint old endpoints.
\item For $D=-12451$ the unique tame line gives an everywhere-finite-unramified cyclic quintic extension of the quartic mirror field. Can a theta-lift or height construction estimate its splitting primes with the sign and square-root accuracy of Section~\ref{sec:joint-test}, first for this line alone? The other five lines have a known conductor at five, so this family distinguishes that first case from its ramified companions.
\end{enumerate}

\begin{thebibliography}{9}
\bibitem{graphEH} P.~L.~Erd\H{o}s, G.~Harcos, S.~R.~Kharel, P.~Maga,
T.~R.~Mezei and Z.~Toroczkai, The sequence of prime gaps is graphic,
\emph{Math. Ann.} \textbf{388} (2024), 2195--2215.
\url{https://doi.org/10.1007/s00208-023-02574-1}.
\bibitem{graphAF} K.~Aggarwal, R.~Frot, H.~Gou and H.~Wang,
Explicit bounds for the graphicality of the prime gap sequence,
arXiv:2512.24230v3 (August 2026).
\url{https://arxiv.org/abs/2512.24230v3}.
\bibitem{graphS} L.~Schoenfeld, Sharper bounds for the Chebyshev functions
$\theta(x)$ and $\psi(x)$. II, \emph{Math. Comp.} \textbf{30} (1976), 337--360.
\url{https://doi.org/10.1090/S0025-5718-1976-0457374-X}.
\bibitem{graphDPG} P.~L.~Erd\H{o}s, S.~R.~Kharel, T.~R.~Mezei and
Z.~Toroczkai, Degree-preserving graph dynamics: a versatile process to
construct random networks, \emph{J. Complex Netw.} \textbf{11} (2023),
cnad046, Section~6. \url{https://doi.org/10.1093/comnet/cnad046}.

\end{thebibliography}
\end{document}
